Latent world models can \emph{decode} physical state from their representations, yet fail to \emph{obey} physical law when predictions are rolled out: a joint-embedding predictive (JEPA-style) model that recovers object position almost perfectly at one step collapses over long horizons and under out-of-distribution (OOD) physics. A popular remedy injects physical inductive biases into the latent space. Testing this remedy on a JEPA world model across five injection mechanisms (fixed physical slots, kinematic dynamics, deep-supervision probes, consistency losses, and label-free priors), two training regimes (post-hoc fine-tuning and from-scratch pretraining), three controlled physics domains, and two photorealistic-simulation benchmarks, we find it fails as a general-purpose lever: it improves prediction in only one of thirty mechanism--domain cells---and there only because it happens to match the domain's dynamics, reversing sign elsewhere---while otherwise \emph{hurting}, with damage growing as the baseline strengthens (up to $+0.56$ normalized MSE from scratch). We trace the failure to a \emph{load-bearing problem}: the black-box channel redundantly encodes the same physical quantities---we verify this directly by decoding position from the non-slot dimensions alone across all three domains---so prediction bypasses the structured slot, and a one-line loss reweighting removes the harm but never yields a net gain. The variables that actually govern physical robustness lie in training and data, not structure: autoregressive free rollout is the single dominant lever ($2.2$--$3.6\times$ OOD error reduction across all domains, seed-robust, and the best transfer to photorealistic scenes), rollout horizons must match dynamics complexity (a $19\times$ long-horizon gain on photorealistic video when paired with geometric scale jitter), and augmentation must match the domain's true nuisances---appearance jitter, the strongest lever on simple synthetic data, \emph{reverses} catastrophically (up to $100\times$) on photorealistic scenes where appearance carries physical information. Finally, we show that cosine- and probe-based metrics are duals of the auxiliary training losses and systematically overstate physical competence, explaining why structural remedies can appear to work.
